\section{Pullback of the original-zeta coefficient sheaf to the complex
projective
line}\label{pullback-of-the-original-zeta-coefficient-sheaf-to-the-complex-projective-line}

Complete derivation, 24 September 2026. Proof locators CSP0--CSP13.

\subsection{CSP0. Input stage, source coverage, and the map being
constructed}\label{csp0.-input-stage-source-coverage-and-the-map-being-constructed}

The support remains \(\tau\langle Z_1;\mathrm{no}\ Z_2\rangle\). This
calculation starts after the complete arithmetic reconstruction, on the
actual coefficient sheaf of CS and CSL. All complex vector spaces,
scalar actions, analytic coordinates and quotient maps below belong to
this receiving construction. They supply no addition, subtraction,
metric, numerical coordinate, or parity on \(\tau\).

The complete corrected argument in
\nolinkurl{USER_ARGUMENT_RECONSTRUCTION.md} and the prerequisite checks
PC01--PC04 govern this use. In particular, fixed support does not imply
a trivial stalk; the retracted addition at \(\tau\) is not used; neither
the counting unit nor selected primes initialize this construction.

The original author source is Alain Connes and Caterina Consani,
\emph{Schemes over \(\mathbb F_1\) and zeta functions},
\href{https://arxiv.org/abs/0903.2024v3}{arXiv:0903.2024v3}, §5. For
this derivation the unchanged author LaTeX file
\nolinkurl{sources/CC_0903_2024_v3/author_source/announc3.tex} was read
at lines1471--1670, including projback, fonction3, restmaps, bord,
hzero, rep1, liftw and court. Its source identity and preserved archive
are recorded in CC\_SHEAF\_SOURCE\_READING\_PRIVATE.json. CS0--CS10 and
CSL0--CSL9 were read in full. These local derivations prove the precise
finite-unit sector, endpoint conventions, and raw Mellin comparison used
here.

The reading of the later Connes--Consani author file CC.tex covered
lines388--423 and485--608: the field-valued point construction, its
three-point support, and the stated antiholomorphic real structure. That
bounded reading is not described as a new reading of the entire paper.
The source is
\href{https://arxiv.org/abs/2609.00299v1}{arXiv:2609.00299v1}. The two
mirrors will remain distinct: the holomorphic \(z\mapsto-1/z\) and the
actual antiholomorphic \(z\mapsto-1/\overline z\) have different
orientation degrees.

The subsequently completed exact-image derivations
\nolinkurl{../quantum_tau_programme_bridge_20260924/EXACT_SCHWARTZ_SUMMATION_IMAGE.md,}
SSI0--SSI10, and \nolinkurl{EXACT_SCHWARTZ_IMAGE_INDEPENDENT_CHECK.md,}
ESI0--ESI12, were read in full. Their proved equality \(J=\overline J\),
continuous inverse, and source-domain qualification are incorporated
below, rather than left as a pending condition.

Let \[
X=\{x_+,\eta,x_-\},\qquad
U_+=\{x_+,\eta\},\quad U_-=\{\eta,x_-\},\quad U=\{\eta\}.
\tag{CSP0.1}
\] Let \(Y=\mathbb P^1(\mathbb C)\) with its usual topology,
\(D=\{0,\infty\}\), and \(Y^\circ=\mathbb C^\times\). Define \[
\pi:Y\longrightarrow X,\qquad
\pi(0)=x_+,\quad\pi(\infty)=x_-,\quad
\pi(z)=\eta\quad(z\ne0,\infty).
\tag{CSP0.2}
\] Its inverse images of the three proper nonempty opens are
\(\mathbb C\), \(Y\setminus\{0\}\), and \(\mathbb C^\times\); hence it
is continuous. It is the evaluation/specialization map for the
projective-line charts: a nonzero field element is a unit and belongs to
their common localization.

In particular, \(\pi^{-1}(\eta)=\mathbb C^\times\). No point of the
complex coordinate plane is identified with \(\tau\), and the complex
chart origin \(0\) maps to \(x_+\), not to \(\eta\). The sphere is a
receiving space over the three-point support, not a replacement
definition of that support.

The object to compute is \[
\mathscr F=\pi^{-1}\Omega
\tag{CSP0.3}
\] in the category of sheaves of complex vector spaces. The coefficient
spaces and every finite complex below retain their specified locally
convex topologies. Derived sheaf operations are taken in that abelian
category; no unproved claim about derived categories of locally convex
spaces is needed.

\subsection{CSP1. Coefficients, endpoint lines, and
restrictions}\label{csp1.-coefficients-endpoint-lines-and-restrictions}

Retain the actual source spaces \[
S=\left\{f\in\mathcal S(\mathbb R;\mathbb C):
f(-v)=f(v),\ f(0)=0,\ \int_{\mathbb R}f(v)\,dv=0\right\},
\] \[
V_+=S\oplus E_+,\quad V_-=S\oplus E_-,
\qquad E_\pm=\mathbb C e_{\pm,0}\oplus\mathbb C e_{\pm,1},
\tag{CSP1.1}
\] and the Fréchet overlap \[
A=\left\{a\in C^\infty(\mathbb R_{>0}):
p_{N,j}(a)=\sup_{u>0}(u^N+u^{-N})
|(u\partial_u)^j a(u)|<\infty,\ N,j\ge0\right\}.
\tag{CSP1.2}
\] The endpoint indices retain value at zero and real integral,
separately on both charts. Put \[
\Sigma f(u)=2\sum_{k\ge1}f(ku),\qquad
Ra(u)=u^{-1}a(u^{-1}),
\] \[
r_+(f,c_0,c_1)=\Sigma f,\qquad
r_-(h,d_0,d_1)=R\Sigma h=\Sigma\widehat h,
\quad
\widehat h(t)=\int_{\mathbb R}h(v)e^{-2\pi ivt}\,dv.
\tag{CSP1.3}
\] The full Poisson formula, before restricting to \(S\), is \[
\Sigma\widehat h(u)
=u^{-1}\Sigma h(u^{-1})+u^{-1}h(0)-\int_{\mathbb R}h(v)\,dv.
\tag{CSP1.4}
\] CS3 proves continuity and injectivity of \(\Sigma:S\to A\); Fourier
transformation preserves \(S\), squares to one there, and gives
\(R\Sigma=\Sigma\widehat{\phantom h}\). Consequently \[
J=\Sigma(S)=r_+(V_+)=r_-(V_-),\qquad
\ker r_\pm=E_\pm.
\tag{CSP1.5}
\] To retain the exact comparison maps used in the earlier source
calculations, write \[
C=\overline J^{\,A},\qquad
Q_{\rm alg}=A/J,\quad Q_H=A/C,\quad N=C/J.
\tag{CSP1.6}
\] The exact row \(0\to N\to Q_{\rm alg}\to Q_H\to0\) has the actual
quotient topologies. The completed SSI1--SSI7 theorem, independently
derived in ESI0--ESI8, now proves \[
\boxed{J=C,\qquad N=0,\qquad Q_{\rm alg}=Q_H=A/J,\qquad
\Sigma:S\xrightarrow{\sim}J\text{ continuously with continuous inverse}.}
\tag{CSP1.7}
\] Here is the exact content being imported, including its factors. For
every entire \(F\) in the full original nontrivial-zero jet ideal,
\(H(s)=F(s)/(2\zeta(s))\) is regular at those zeros and at \(1\), and
has possible residues \(F(-2r)/(2\zeta'(-2r))\) at the trivial zeros.
SSI4 proves rapid vertical bounds for \(P_N(s)H(s)\),
\(P_N(s)=\prod_{r=1}^N(s+2r)\), by the full functional equation and its
subquadratic-strip maximum-principle estimate. SSI5--SSI6 then construct
the actual Schwartz inverse \[
f_F(x)=\frac1{2\pi}\int_{\mathbb R}
\frac{F(2+it)}{2\zeta(2+it)}x^{-2-it}\,dt,\quad x>0,
\qquad f_F(-x)=f_F(x),
\] \[
\frac{f_F^{(2r)}(0)}{(2r)!}
=\frac{F(-2r)}{2\zeta'(-2r)},\quad
f_F(0)=0,\quad \int_{\mathbb R}f_F=0,\quad
\mathcal M_0\Sigma f_F=F.
\tag{CSP1.8}
\] The derivative remainder and finite-seminorm bounds there prove that
this is a continuous inverse on precisely \(S\), not an approximation in
a closure. Equivalently, for \(a\in J\) and \(x>0\),
\(\Sigma^{-1}a(x)=\tfrac12\sum_{k\ge1}\mu(k)a(kx)\), with the
zero-endpoint extension supplied by (CSP1.8), not by termwise
evaluation. Those complete proofs supply (CSP1.7). Ordinary cohomology
below is therefore already Hausdorff; retaining both earlier quotient
names does not leave an unresolved extra kernel.

The sheaf \(\Omega\) is the diagram
\(V_+\xrightarrow{r_+}A\xleftarrow{r_-}V_-\), with global sections the
fibre product. CS0--CS2A specify the full-difference interpretation of
the original adelic coinvariants used to identify this source sector.
Under the alternative moment-null generating-domain reading, the extra
kernel computed there must be retained; this note uses the stated CS
source, not that alternative.

\subsection{CSP2. An explicit construction of the inverse-image
sheaf}\label{csp2.-an-explicit-construction-of-the-inverse-image-sheaf}

Write \(j:Y^\circ\hookrightarrow Y\), \(i:D\hookrightarrow Y\), and
\(\underline A_Y,\underline A_{Y^\circ}\) for constant sheaves, whose
sections are locally constant functions with values in the vector space
\(A\). The coefficient topology does not change the meaning of ``locally
constant.''

The natural map \[
\underline A_Y\longrightarrow j_*\underline A_{Y^\circ}
\tag{CSP2.1}
\] is an isomorphism. Away from the poles this is immediate. At either
pole, arbitrarily small punctured disks are connected, so the right-hand
stalk is \(A\), and the map of stalks is the identity. A stalkwise
isomorphism of sheaves is an isomorphism. This is an assertion about
\(j_*\), not an assertion that \(Rj_*\) has no higher terms.

Define \(i_*(V_+\oplus V_-)\) as the direct sum of the two skyscraper
sheaves, and let \[
\epsilon:\underline A_Y\longrightarrow i_*(A\oplus A)
\] be restriction to the two pole stalks. Then \[
\boxed{\mathscr F
=\underline A_Y\times_{\,i_*(A\oplus A)}i_*(V_+\oplus V_-),}
\qquad
(a,v_+,v_-)\mapsto
\bigl(a|_0-r_+v_+,\,a|_\infty-r_-v_-\bigr).
\tag{CSP2.2}
\] Here and below a pole component is omitted on an open set not
containing that pole.

To prove that (CSP2.2) is (CSP0.3), first construct a map
\(\Omega\to\pi_*\mathscr F\). A section \(v_+\) on \(U_+\) gives the
locally constant section \(r_+v_+\) on \(\pi^{-1}U_+\), together with
pole value \(v_+\). Use \(r_-v_-\) on the other chart and the identity
on their overlap. The fibre-product condition makes these maps
compatible, so adjunction gives \(\pi^{-1}\Omega\to\mathscr F\). Its
pole stalks are the identities on \(V_\pm\), and its other stalks are
the identity on \(A\). It is therefore an isomorphism.

In particular, on a connected disk containing only one pole, sections
are exactly \(V_\pm\), with their original Fréchet topology. On a
connected open avoiding both poles, sections are \(A\). The restrictions
into the punctured disk are exactly \(r_\pm\), not injections invented
from a coordinate picture.

Formula (CSP2.2) gives a short exact sequence of vector-space sheaves \[
0\longrightarrow\mathscr F
\longrightarrow\underline A_Y\oplus i_*(V_+\oplus V_-)
\xrightarrow{\ \epsilon-r\ }i_*(A\oplus A)
\longrightarrow0.
\tag{CSP2.3}
\] Surjectivity is stalkwise: at each pole the constant-sheaf
\(A\)-coordinate supplies any prescribed value, independently of the
restriction map \(r_\pm\). All displayed local maps are continuous. This
construction does not assume that \(J\) is closed or that
\(\Sigma^{-1}\) is continuous.

\subsection{CSP3. Disk, annulus, and sphere cohomology with these
coefficients}\label{csp3.-disk-annulus-and-sphere-cohomology-with-these-coefficients}

Here is the local resolution used in the calculation. For a complex
vector space \(B\), let \(\mathscr E_B^q\) be the sheaf of differential
forms locally expressible as a finite sum \(\sum_\nu\omega_\nu b_\nu\),
with ordinary smooth complex forms \(\omega_\nu\) and vectors
\(b_\nu\in B\). Its differential acts on the forms. On a star-shaped
coordinate ball, the radial homotopy \[
(K\omega)_x=\int_0^1 t^{q-1}
\,\iota_x\omega_{tx}\,dt\quad(q>0)
\tag{CSP3.1}
\] acts on those finitely many scalar coefficients and satisfies
\(dK+Kd=1\) in positive degree and \(Kd=1-\operatorname{ev}_0\) in
degree zero. Thus \[
0\to\underline B\to\mathscr E_B^0\xrightarrow d
\mathscr E_B^1\xrightarrow d\cdots
\tag{CSP3.2}
\] is a resolution. Partitions of unity act by multiplication, so these
sheaves are fine. A locally finite partition produces locally finite
sums, which are legitimate sections of the sheaf. Hence this computes
sheaf cohomology. It also applies to a possibly noncomplete subspace
\(J\) or quotient \(Q_{\rm alg}\): only finitely many coefficient
vectors occur locally.

For the complete Fréchet spaces \(A,V_\pm\), one may equivalently use
smooth Fréchet-valued forms. The same integral exists as a limit of
Riemann sums in every seminorm, and \[
p\!\left(\int_a^b f(t)\,dt\right)
\le\int_a^b p(f(t))\,dt
\tag{CSP3.3}
\] proves continuity of each contraction on compact sets. Finite good
covers and their contractions consequently give the same finite
coefficient complexes with continuous maps. No completed infinite tensor
product is being substituted for a coefficient space.

A disk, or a plane chart with its pole included, is acyclic for a
constant sheaf: (CSP3.1) contracts its positive-degree complex. On
\(Y^\circ\), put \(z=e^{t+i\theta}\). A closed one-form can be written
on the covering \(\mathbb R^2\) as \[
\omega=P(t,\theta)\,dt+Q(t,\theta)\,d\theta,
\quad
\partial_\theta P=\partial_tQ,
\] with \(2\pi\)-periodic coefficients. Its period \[
a=\int_0^{2\pi}Q(0,\theta)\,d\theta
\tag{CSP3.4}
\] is independent of \(t\). After subtracting \(a\,d\theta/(2\pi)\), the
explicit primitive \[
g(t,\theta)=
\int_0^tP(x,0)\,dx+
\int_0^\theta\left(Q(t,\phi)-a/(2\pi)\right)d\phi
\tag{CSP3.5}
\] is periodic and has differential equal to the adjusted form. Every
two-form \(G(t,\theta)\,dt\wedge d\theta\) is the differential of
\((\int_0^tG(x,\theta)dx)d\theta\). These formulas prove \[
R\Gamma(Y^\circ,\underline B)
\simeq B[0]\oplus B[-1].
\tag{CSP3.6}
\] The degree-one coordinate is the positive \(z\)-circle period
(CSP3.4), with representative \(a\,d\theta/(2\pi)\). They also prove
continuity of the period and chosen primitives in the coefficient
seminorms on compact sets. A punctured pole disk has the same
calculation in its local coordinate.

For the sphere, cover by its two plane charts and use (CSP3.2) to form
the Čech double complex. Each chart is acyclic for \(\underline B\);
their overlap has (CSP3.6). The constant restriction difference is
\(B\oplus B\to B,\ (a,b)\mapsto a-b\), which is onto. Consequently \[
H^0(Y,\underline B)=B,\qquad
H^1(Y,\underline B)=0,\qquad
H^2(Y,\underline B)=B,\qquad H^{>2}=0.
\tag{CSP3.7}
\] The degree-two coordinate is integration with the complex
orientation. For a concrete generator one can use \[
\omega_Y=\frac1\pi\frac{dx\wedge dy}{(1+|z|^2)^2},
\quad z=x+iy,\qquad \int_Y\omega_Y=1.
\tag{CSP3.8}
\] It extends smoothly through infinity. The integral is
\(2\int_0^\infty r(1+r^2)^{-2}dr=1\). The preceding two-chart complex
shows that its nonzero integral generates the entire degree-two group; a
zero-integral form is exact by that same period calculation. Thus
(CSP3.7) is established by the resolution and its maps, not inferred
solely from a homotopy picture.

\subsection{CSP4. Global cohomology and a continuous finite
complex}\label{csp4.-global-cohomology-and-a-continuous-finite-complex}

A skyscraper sheaf on either pole has no higher cohomology: sections on
an open are either its coefficient space or zero, and all restrictions
are onto. Apply (CSP3.7) to (CSP2.3). A finite model for
\(R\Gamma(Y,\mathscr F)\) is \[
K^0=A\oplus V_+\oplus V_-,\qquad
K^1=A\oplus A,\qquad K^2=A,
\] \[
d_K^0(a,v_+,v_-)=
(a-r_+v_+,\,a-r_-v_-),\qquad d_K^1=0.
\tag{CSP4.1}
\] The degree-two coordinate comes from the constant-sheaf summand with
the orientation (CSP3.8). There are no other degrees.

The following explicit contraction removes the redundant diagonal \(A\).
Define \[
D^0=V_+\oplus V_-,\qquad D^1=A,\qquad D^2=A,
\qquad d_D^0=r_+-r_-,\quad d_D^1=0.
\tag{CSP4.2}
\] Maps \(p:K\to D\) and \(s:D\to K\) are \[
p^0(a,v_+,v_-)=(v_+,v_-),\quad
p^1(b_+,b_-)=b_--b_+,\quad p^2=\mathrm{id},
\] \[
s^0(v_+,v_-)=(r_-v_-,v_+,v_-),\quad
s^1(b)=(-b,0),\quad s^2=\mathrm{id}.
\tag{CSP4.3}
\] They are chain maps and \(ps=1\). Let \(h^1(b_+,b_-)=(b_-,0,0)\),
with every other component of \(h\) zero. Direct substitution gives \[
1-sp=d_Kh+hd_K.
\tag{CSP4.4}
\] Every map is continuous and uses only the original restrictions,
additions in coefficient spaces, and finite products. No inverse to
\(\Sigma\) is needed.

It follows that \[
\boxed{
\begin{aligned}
H^0(Y,\mathscr F)&=
\{((\widehat h,c_0,c_1),(h,d_0,d_1)):
h\in S,\ c_0,c_1,d_0,d_1\in\mathbb C\},\\
H^1(Y,\mathscr F)&=A/J=Q_{\rm alg},\\
H^2(Y,\mathscr F)&=A,\qquad H^{>2}(Y,\mathscr F)=0.
\end{aligned}}
\tag{CSP4.5}
\] The first isomorphism is topological via the Fourier map and
coordinate projections. Degree one has its original quotient topology,
which is Fréchet and Hausdorff by (CSP1.7); the previously retained
comparison subspace \(C/J\) is zero. Degree two has the original Fréchet
topology of \(A\), and its isomorphism is induced by the sheaf
projection \(\mathscr F\to\underline A_Y\), followed by complex-oriented
integration.

Equivalently, two pole charts for \(\mathscr F\) are acyclic: on a
chart, (CSP2.3) reduces to the complex
\(A\oplus V_\pm\to A,\ (a,v)\mapsto a-r_\pm v\), contracted onto
\(V_\pm\). Their annular overlap has \(A[0]\oplus A[-1]\). In the
derived Mayer--Vietoris cone its degree-zero overlap produces \(A/J\),
and its degree-one loop produces the additional \(A\) in degree two.
This is the exact Čech-loop contribution.

\subsection{CSP5. Where the additional term lives over the original
support}\label{csp5.-where-the-additional-term-lives-over-the-original-support}

Because each point of \(X\) has a smallest open neighborhood, higher
direct-image stalks are the cohomology of the corresponding preimages.
The preceding chart and overlap calculations give \[
R^0\pi_*\mathscr F=\Omega,\qquad
R^1\pi_*\mathscr F=j_{X!}\underline A_{\{\eta\}},
\qquad R^{q>1}\pi_*\mathscr F=0,
\tag{CSP5.1}
\] where \(j_X:\{\eta\}\hookrightarrow X\). Indeed the degree-one sheaf
has zero stalks at \(x_\pm\) and stalk \(A\) at \(\eta\), with the zero
generization maps from the closed stalks.

Its Čech complex on \(X\) is simply \(0\to A\), in degrees zero and one,
so \[
H^0(X,R^1\pi_*\mathscr F)=0,\qquad
H^1(X,R^1\pi_*\mathscr F)=A.
\tag{CSP5.2}
\] The two-column, two-row Čech resolution computing the composite of
direct image and global sections now has no possible nonzero
differential beyond its displayed maps. It gives the Leray
identifications \[
H^0(Y,\mathscr F)=H^0(X,\Omega),\quad
H^1(Y,\mathscr F)=H^1(X,\Omega),\quad
H^2(Y,\mathscr F)=H^1(X,R^1\pi_*\mathscr F)=A.
\tag{CSP5.3}
\] The first two maps are induced by the adjunction unit
\(\Omega\to R\pi_*\pi^{-1}\Omega\). This supplies an exact derived
comparison with the finite-poset calculation: the extra term is the
punctured-disk loop in the actual receiver. It is not inserted into the
stalk of \(\tau\) as an additional parity label.

For comparison, \(R^0j_*\underline A_{Y^\circ}=\underline A_Y\), whereas
\(R^1j_*\underline A_{Y^\circ}\) has one \(A\)-stalk at each pole.
Formula (CSP2.1) did not discard those higher direct images. No
splitting of the whole object \(Rj_*\underline A\) is inferred from the
list of its cohomology sheaves.

\subsection{CSP6. Exact pole costalks and their orientation
convention}\label{csp6.-exact-pole-costalks-and-their-orientation-convention}

Let \(p\) be either pole, \(B_p\) a small disk with local coordinate
\(z\) at \(0\), or \(w=1/z\) at \(\infty\). Use the positive
local-coordinate circle in each case. Local cohomology is the fibre
complex of restriction: \[
R\Gamma_{\{p\}}(B_p,\mathscr F)
=\operatorname{Cone}\bigl(R\Gamma(B_p,\mathscr F)
\to R\Gamma(B_p\setminus\{p\},\mathscr F)\bigr)[-1].
\tag{CSP6.1}
\] Using the explicit resolutions above, choose the fibre-complex
convention whose differential is the restriction map. It gives \[
L_p^0=V_p,\qquad L_p^1=A,\qquad L_p^2=A,
\qquad d^0=r_p,\quad d^1=0.
\tag{CSP6.2}
\] In the unreoriented fibre cone the connecting map inserts the
punctured-disk coefficient with a positive sign in the next degree. We
keep that coordinate in degree one. In degree two we take its negative,
so that the coordinate is positive complex-oriented integration, as
checked next. The zero differential into degree two is unchanged. Thus
\[
H^0_{\{p\}}=E_p,\qquad
H^1_{\{p\}}=Q_{\rm alg},\qquad
H^2_{\{p\}}=A,\qquad H^{>2}_{\{p\}}=0.
\tag{CSP6.3}
\] These compute the costalk \(i_p^!\mathscr F\); passage to smaller
pole disks gives identity maps on these models.

The degree-two coordinate is the positive complex local fundamental
class. To check the sign of the connecting map, choose a radial function
\(\chi\) that is one near the pole and zero near the outside of a small
disk. If \(\theta_p\) is its positive local angle, then \[
-d\chi\wedge\frac{d\theta_p}{2\pi}
\tag{CSP6.4}
\] is a compactly supported smooth two-form away from the pole. Its
integral is \(-\int_0^\infty\chi'(r)dr=1\). Put
\(\beta=(1-\chi)d\theta_p/(2\pi)\), a smooth one-form on the disk, so
\(d\beta\) is (CSP6.4). Restrict the relative cone to an outer annular
collar, where \(\beta=d\theta_p/(2\pi)\). Its cone differential sends
\((\beta,0)\) to \((d\beta,d\theta_p/(2\pi))\). Hence the boundary class
\((0,d\theta_p/(2\pi))\) equals \((-d\beta,0)\) in that relative
cohomology. It is the \textbf{negative} of the positive-integral class.
Excision identifies the same sign in the point-support group.

Thus the local connecting map from a positive angular period \(a\) is
\(-a\) in our degree-two integration coordinate. The support-to-global
map on that coordinate is \(+1\) at both poles, because both positive
local complex orientations agree with the sphere orientation. At
infinity the local circle is the positive \(w\)-circle, not the positive
\(z\)-circle.

\subsection{CSP7. Both supports, the complete localization triangle, and
all
signs}\label{csp7.-both-supports-the-complete-localization-triangle-and-all-signs}

For \(D=\{0,\infty\}\), take the direct sum of the disjoint costalk
models: \[
L_D^0=V_+\oplus V_-,\quad
L_D^1=A\oplus A,\quad L_D^2=A\oplus A,
\] \[
d_L^0(v_+,v_-)=(r_+v_+,r_-v_-),\qquad d_L^1=0.
\tag{CSP7.1}
\] The support-to-global map \(k:L_D\to D^\bullet\) from (CSP4.2) is \[
k^0=\mathrm{id},\qquad
k^1(b_+,b_-)=b_+-b_-,\qquad
k^2(c_+,c_-)=c_++c_-.
\tag{CSP7.2}
\] Degree one has the ordered-chart difference; degree two has the sum
of the two positive local orientation classes. These are different signs
in different degrees.

An explicit verification of the entire triangle is useful. The ordinary
cone of \(k\) has \[
\operatorname{Cone}(k)^{-1}=V_+\oplus V_-,
\] \[
\operatorname{Cone}(k)^0=(V_+\oplus V_-)\oplus A^2,\quad
\operatorname{Cone}(k)^1=A\oplus A^2,\quad
\operatorname{Cone}(k)^2=A,
\] with \[
d^{-1}(v)=(v,-r_+v_+,-r_-v_-),
\] \[
d^0(v,b_+,b_-)=(r_+v_+-r_-v_-+b_+-b_-,\,0,0),
\quad
d^1(b,c_+,c_-)=c_++c_-.
\tag{CSP7.3}
\] The overlap model \(A[0]\oplus A[-1]\) maps into this cone by \[
a\ \text{in degree0}\longmapsto (0,a,a),\qquad
a\ \text{in degree1}\longmapsto (0,-a,a).
\tag{CSP7.4}
\] These are cocycles. Every degree-zero cocycle differs by
\(d^{-1}(v)\) from a unique diagonal pair \((0,a,a)\). Every degree-one
cocycle has \(c_-=-c_+\); its \(b\)-component is removed by the boundary
of \((0,b,0)\) in degree zero. Degrees minus one and two have zero
cohomology by injectivity of \(d^{-1}\) and surjectivity of the sum.
Thus (CSP7.4) is a quasi-isomorphism, with all reductions given by
continuous finite linear operations.

This proves the localization triangle and its boundaries in the chosen
coordinates. In particular, \[
\partial^0:A\to Q_{\rm alg}^2,\quad a\mapsto([a],[a]),
\] \[
\partial^1:H^1(Y^\circ,\underline A)=A
\to H^2_D(Y,\mathscr F)=A^2,\quad a\mapsto(-a,a).
\tag{CSP7.5}
\] The second formula combines the negative localization-boundary sign
proved in CSP6 with \(w=1/z\), which gives \(d\arg w=-d\arg z\).

The full nonzero portion of the ordinary localization sequence is \[
\begin{aligned}
0\to E_+\oplus E_-&\to H^0(Y,\mathscr F)
 \xrightarrow{\ r\ }A
 \xrightarrow{\ \mathrm{diag}\circ q\ }Q_{\rm alg}^2\\
&\xrightarrow{\ (q_+,q_-)\mapsto q_+-q_-\ }Q_{\rm alg}
 \xrightarrow{\ 0\ }A
 \xrightarrow{\ a\mapsto(-a,a)\ }A^2
 \xrightarrow{\ (a,b)\mapsto a+b\ }A\to0.
\end{aligned}
\tag{CSP7.6}
\] Here \(r\) takes a Fourier-graph section to its common restriction
and has image \(J\); \(q:A\to A/J\). Every kernel and image in this row
follows directly from the displayed formulas. In particular, global
degree-one classes restrict to zero in annular degree one: they are Čech
differences of degree-zero overlap data.

For one pole alone the support-to-global maps are \[
L_+\to D:\quad(v,b,c)\mapsto((v,0),b,c),
\] \[
L_-\to D:\quad(v,b,c)\mapsto((0,v),-b,c).
\tag{CSP7.7}
\] They induce degree-one isomorphisms with signs \(+1,-1\) and
degree-two isomorphisms with sign \(+1\) in both cases. The
complementary chart is acyclic for \(\mathscr F\), as required by its
localization sequence.

\subsection{CSP8. Holomorphic power maps and the geometric normal
factor}\label{csp8.-holomorphic-power-maps-and-the-geometric-normal-factor}

For each recovered positive integer \(n\), the actual holomorphic map \[
b_n:Y\to Y,\qquad b_n(z)=z^n,\quad b_n(0)=0,\quad b_n(\infty)=\infty
\tag{CSP8.1}
\] satisfies \(\pi b_n=\pi\), so \(b_n^{-1}\mathscr F\) is canonically
identified with \(\mathscr F\) before applying any coefficient action.

For coefficients retain \[
T_a f(u)=f(u/a),\qquad
\rho_+(a)(f,c_0,c_1)=(f(\cdot/a),c_0,a c_1),
\] \[
\rho_-(a)(h,d_0,d_1)=(a h(a\,\cdot),a d_0,d_1),\qquad a>0.
\tag{CSP8.2}
\] CS8 proves \(r_\pm\rho_\pm(a)=T_ar_\pm\). These give a continuous
sheaf endomorphism over the identity on \(Y\). Combine it with pullback
by \(b_n\), using \(a=n\), and denote the resulting cohomological
operator by \(\mathsf B_n\).

On an annular angular class, \[
b_n^*\left(a\,\frac{d\theta}{2\pi}\right)
=n\,a\,\frac{d\theta}{2\pi}.
\tag{CSP8.3}
\] The local coordinate at infinity also maps as \(w\mapsto w^n\), so
the same positive factor holds for its positive local circle. On degree
two, either (CSP6.4) and the boundary give factor \(n\), or one can
compute its global degree directly: \[
b_n^*\omega_Y
=\frac1\pi
\frac{n^2r^{2n-2}\,dx\wedge dy}{(1+r^{2n})^2},\qquad
\int_Yb_n^*\omega_Y
=2n^2\int_0^\infty\frac{r^{2n-1}}{(1+r^{2n})^2}dr=n.
\tag{CSP8.4}
\] All constants and the branch multiplicities at both poles are
included. Therefore the induced chain/cohomology model for the combined
operation is \[
\mathsf B_n^0=\rho_+(n)\oplus\rho_-(n),\qquad
\mathsf B_n^1=T_n,\qquad
\mathsf B_n^2=nT_n.
\tag{CSP8.5}
\] Each pole-support complex has the corresponding actions
\((\rho_\pm(n),T_n,nT_n)\). On the overlap, degree zero has \(T_n\) and
degree one has \(nT_n\); these make every map in (CSP7.6) equivariant.

This proves, as a representation of the recovered positive-integer
monoid, \[
\boxed{H^2(Y,\mathscr F)=A(-1),\qquad
A(-1)\ \text{has action }nT_n.}
\tag{CSP8.6}
\] The notation records the proved degree character, not a declaration
that this infinite-dimensional complex sheaf is an étale Tate sheaf. The
coefficient-only action remains \(T_a\) in every degree. The additional
\(n\) in (CSP8.6) is produced by the actual geometric covering, not by
changing the coefficient action by definition.

Composition is exact: \(b_m b_n=b_{mn}\), \(T_mT_n=T_{mn}\), and hence
\(\mathsf B_m\mathsf B_n=\mathsf B_{mn}\). The induced operators are
invertible on all displayed cohomology spaces. On positive rational
indices define the cohomological comparison
\(\mathsf B_{n/m}=\mathsf B_n\mathsf B_m^{-1}\); equality \(nm'=n'm\)
proves independence of the presentation. In degree two its action is
\((n/m)T_{n/m}\). The strongly continuous family \(aT_a\), already
defined on the coefficient space, is the unique extension of this
rational action to \(a>0\). This does not assert a holomorphic map
\(z\mapsto z^a\) for a nonintegral real \(a\).

The maps \(b_n\) here are receiving-space covers. The full signed CC
structure has its own odd-index and chart-swapping rules; even \(b_n\),
or all \(b_n\) with one uniform chart rule, are not thereby proved to be
endomorphisms of that full signed structure.

\subsection{CSP9. The two mirrors and their different orientation
signs}\label{csp9.-the-two-mirrors-and-their-different-orientation-signs}

Keep the maps separate: \[
w(z)=-1/z,\qquad
\alpha(z)=-1/\overline z.
\tag{CSP9.1}
\] Both exchange the two poles and induce the same swap
\(x_+\leftrightarrow x_-\) on \(X\). The former is holomorphic, degree
\(+1\) on the sphere; the latter is antiholomorphic, degree \(-1\), and
is the actual real structure stated in CC.tex. On the unit circle their
angle maps are respectively \[
\theta\mapsto\pi-\theta,\qquad \theta\mapsto\pi+\theta.
\tag{CSP9.2}
\] Thus their actions on \(H^1(Y^\circ,\underline A)\), before a
coefficient map, have signs \(-1,+1\), respectively.

Use the actual source Weyl coefficient map: exchange \(V_+\) and \(V_-\)
by identity on their correspondingly labelled coordinates, and apply
\(R\) on \(A\). It is a sheaf map over either swap because \(r_-=Rr_+\)
and \(R^2=1\). It is complex linear. It is not being identified without
further construction with the full signed structure-sheaf action. If a
semilinear coefficient comparison is wanted, composing with complex
conjugation gives a separately specified conjugate-linear map.

On the global model (CSP4.2) the two operations are \[
\begin{array}{c|ccc}
 &D^0&D^1&D^2\\ \hline
w&\text{swap}&-R&+R\\
\alpha&\text{swap}&-R&-R.
\end{array}
\tag{CSP9.3}
\] The degree-one sign is the ordered Čech swap in both cases. Degree
two is the product of that Čech sign with the angular sign in (CSP9.2).
Equivalently it is the orientation degree of the map, times \(R\). On
the pole-support degree-two pair, \(w\) acts as
\((c_+,c_-)\mapsto(Rc_-,Rc_+)\), whereas \(\alpha\) acts as
\((c_+,c_-)\mapsto(-Rc_-,-Rc_+)\). Both act on its degree-one pair by
\((b_+,b_-)\mapsto(Rb_-,Rb_+)\).

These formulas respect all localization maps. For example, \[
(Rb_--Rb_+)=-R(b_+-b_-).
\] For the annular boundary \(a\mapsto(-a,a)\), the holomorphic
operation sends its value to \((Ra,-Ra)\), which is the boundary of
\(-Ra\); the antiholomorphic operation sends it to \((-Ra,Ra)\), the
boundary of \(+Ra\). The degree-two sum is respectively transformed by
\(+R\) and \(-R\).

Thus equality of the finite-poset swap does not erase the sphere's
complex-orientation information. In particular, using \(+R\) on global
degree two for the actual antiholomorphic \(\alpha\) would give the
wrong map.

\subsection{CSP10. The exact quotient of the normal
direction}\label{csp10.-the-exact-quotient-of-the-normal-direction}

For any vector quotient \(B=A/J\), there is an actual sheaf morphism \[
q_Y:\mathscr F\longrightarrow j_!\underline{Q_{\rm alg}}_{Y^\circ}.
\tag{CSP10.1}
\] On \(Y^\circ\) it is \(q:A\to A/J\); at either pole its stalk map is
zero. To verify that it is a sheaf morphism into extension by zero, take
a section near a pole. Its nearby constant \(A\)-value is
\(r_\pm v\in J\), so its quotient vanishes on a neighborhood of that
pole. Thus the quotient section has exactly the vanishing-near-boundary
condition required by \(j_!\). The map is onto on every stalk.

Its kernel \(\mathscr F_J\) has pole stalks \(V_\pm\), interior stalk
\(J\), and the same restrictions with codomain \(J\). Hence \[
0\to\mathscr F_J\to\mathscr F
\xrightarrow{q_Y}j_!\underline{Q_{\rm alg}}\to0
\tag{CSP10.2}
\] is exact. Give \(J\) its actual closed-subspace Fréchet topology from
(CSP1.7). The restrictions \(V_\pm\to J\) are continuous, with
continuous right inverses obtained from \(\Sigma^{-1}\) and Fourier
transformation. CSP3 therefore applies with these complete coefficient
spaces as well as in the vector-space category.

The extension-by-zero sheaf has the same pole/interior model with zero
pole spaces and overlap \(Q_{\rm alg}\). Thus \[
H^0(Y,j_!\underline{Q_{\rm alg}})=0,\quad
H^1(Y,j_!\underline{Q_{\rm alg}})=Q_{\rm alg},\quad
H^2(Y,j_!\underline{Q_{\rm alg}})=Q_{\rm alg}.
\tag{CSP10.3}
\] For \(\mathscr F_J\), both restrictions are onto \(J\); therefore its
degree one is zero and its degree two is \(J\). Under these models
(CSP10.1) induces the identity \(Q_{\rm alg}\to Q_{\rm alg}\) in degree
one and \(q:A\to A/J\) in degree two. In particular the normal-direction
sequence is exactly \[
\boxed{
0\longrightarrow J(-1)\longrightarrow A(-1)
\longrightarrow H^2(Y,j_!\underline{Q_{\rm alg}})
=Q_{\rm alg}(-1)\longrightarrow0.
}
\tag{CSP10.4}
\] Its kernel is the full original summation image \(J\), not a selected
family of zero jets or a finite spectral truncation. The sphere
orientation and the covers prove the \((-1)\) action on each term. As
locally convex spaces, the quotient topology in (CSP10.4) is precisely
the topology of \(A/J\).

There is also the continuous Hausdorff version \[
0\longrightarrow C(-1)\longrightarrow A(-1)
\longrightarrow Q_H(-1)\longrightarrow0,
\tag{CSP10.5}
\] and an exact comparison \[
0\longrightarrow (C/J)(-1)\longrightarrow Q_{\rm alg}(-1)
\longrightarrow Q_H(-1)\longrightarrow0.
\tag{CSP10.6}
\] The sheaf map in this version is the composition with
\(Q_{\rm alg}\to Q_H\); it does not rely on any general exactness
assertion for Hausdorffification. By (CSP1.7), this last map is the
identity comparison, \(C/J=0\), and (CSP10.5) is exactly (CSP10.4). Both
retained kernel labels are the same actual closed source image
\(J(-1)\); there is no remaining completion operation.

\subsection{CSP11. Original zeta, all factors, and the map to the
shifted
kernel}\label{csp11.-original-zeta-all-factors-and-the-map-to-the-shifted-kernel}

Use the exact raw transform and its CS comparison \[
\mathcal M_0a(s)=\int_0^\infty a(u)u^s\,\frac{du}{u},
\qquad \Theta a(s)=\tfrac12\mathcal M_0a(s).
\tag{CSP11.1}
\] The factor \(1/2\) is the explicitly proved comparison map; the
original restriction remains \(\Sigma f=2\sum_{k\ge1}f(ku)\). It has not
been replaced by a different restriction. The target is \[
\mathcal B=\left\{F\text{ entire}:
\sup_{|\Re s|\le M}(1+|\Im s|)^N|F(s)|<\infty
\text{ for every }M,N\ge0\right\}.
\tag{CSP11.2}
\] CS7 and S1 give a topological isomorphism \(\Theta:A\to\mathcal B\),
with inverse \[
(\Theta^{-1}F)(u)=\frac{u^{-1/2}}{\pi}
\int_{\mathbb R}F(1/2+it)u^{-it}\,dt.
\tag{CSP11.3}
\] For \(\Re s>1\), the original identity is \[
\mathcal M_0\Sigma f(s)
=2\zeta(s)\int_0^\infty f(v)v^s\,\frac{dv}{v},
\qquad
\Theta\Sigma f(s)=\zeta(s)\int_0^\infty f(v)v^s\,\frac{dv}{v}.
\tag{CSP11.4}
\] It retains the original Euler product, unit summand and every
repetition: \[
\zeta(s)=1+\sum_{k\ge2}k^{-s}=\prod_p(1-p^{-s})^{-1},
\quad
-\frac{\zeta'(s)}{\zeta(s)}
=\sum_p\sum_{m\ge1}(\log p)p^{-ms}
\quad(\Re s>1).
\tag{CSP11.5}
\]

For the actual retained source \[
f_0(v)=\frac{\pi}{2}v^2(2\pi v^2-3)e^{-\pi v^2},
\] \[
\Theta\Sigma f_0(s)=F_0(s)
=\frac{s(s-1)}8\pi^{-s/2}\Gamma(s/2)\zeta(s),
\qquad \mathcal M_0\Sigma f_0=2F_0.
\tag{CSP11.6}
\] The original zeta pole at \(1\), its trivial zeros, and the Gamma
poles give \[
F_0(0)=F_0(1)=\frac18,
\] \[
F_0(-2r)=
\frac{r(2r+1)(-1)^r\pi^r}{2\,r!}\zeta'(-2r)
=\frac{(1+2r)(2r)}8\pi^{-(1+2r)/2}
\Gamma((1+2r)/2)\zeta(1+2r)\ne0
\quad(r\ge1).
\tag{CSP11.7}
\] For general \(f\in S\), writing
\(a_f(s)=\int_0^\infty f(v)v^s\,dv/v\) with its Taylor-subtracted
continuation, the exact receiver values are \[
\mathcal M_0\Sigma f(-2r)
=2\zeta'(-2r)\frac{f^{(2r)}(0)}{(2r)!},\quad
\mathcal M_0\Sigma f(1)=2a_f'(1),\quad
\mathcal M_0\Sigma f(0)=-a_f(0).
\tag{CSP11.8}
\] These identities keep the endpoint and trivial-zero contributions
rather than representing them as absent zeros of a replacement function.

Let \(\mathscr Z\) be the actual nontrivial zero set of the original
\(\zeta\), with full multiplicities \(m_\rho\), and put \[
\mathcal I=\{F\in\mathcal B:
F^{(j)}(\rho)=0\text{ for every }\rho\in\mathscr Z,\
0\le j<m_\rho\},\qquad
\mathcal Q=\mathcal B/\mathcal I.
\tag{CSP11.9}
\] The exact-image theorem SSI1--SSI7 strengthens S5 and CS7 to
\(\Theta J=\Theta C=\mathcal I\). The scalar \(1/2\) in \(\Theta\)
leaves this linear subspace unchanged. Hence \[
Q_H\xrightarrow{\ \sim\ }\mathcal Q,\quad
[a]\longmapsto[\Theta a],\qquad
\Theta T_a=a^s\Theta.
\tag{CSP11.10}
\] The normal quotient from the actual sphere therefore has the
continuous map \[
\nu:A(-1)\longrightarrow\mathcal Q(-1),\qquad
a\longmapsto[\Theta a].
\tag{CSP11.11}
\] Its kernel is exactly \(J(-1)=C(-1)\) by (CSP1.7). Thus its factor
through the ordinary geometric quotient (CSP10.4) is a continuous
isomorphism \(Q_{\rm alg}(-1)\to\mathcal Q(-1)\). The earlier comparison
kernel \((C/J)(-1)\) is zero by the actual inverse calculation, not by
an assumed exactness property of Hausdorffification.

The comparison with the already constructed high-weight shifted quotient
in GSL6.5--GSL6.11 is now an actual map. Retain \[
\mathcal I_+=\{F\in\mathcal B:
F^{(j)}(\rho+1)=0,\ 0\le j<m_\rho\},\quad
\mathcal Q_{\cup,+}=\mathcal B/(\mathcal I\cap\mathcal I_+).
\] The GSL multiplier \(E_+\) has jet one on the original divisor and
jet zero on the divisor shifted by \(+1\). Its construction and global
strip estimates, not just these jet conditions, are the input
guaranteeing a continuous multiplier of \(\mathcal B\). Define \[
\boxed{
\mathfrak n:A(-1)\longrightarrow\mathcal Q_{\cup,+},\qquad
\mathfrak n(a)=
[(1-E_+(s))(\Theta a)(s-1)]_{\mathcal I\cap\mathcal I_+}.
}
\tag{CSP11.12}
\] This is exactly \(j_{\rm high}\circ\nu\), where GSL proves that
\(j_{\rm high}\) is injective. Thus its kernel is the full original
source \(J(-1)=C(-1)\), its image is the entire GSL kernel, and its
induced map from the ordinary geometric quotient (CSP10.4) is an
isomorphism onto that kernel. Its continuity follows from (CSP11.3),
translation continuity and the stated multiplier bounds.

For every recovered \(n\), equivariance includes the full degree factor:
\[
(\Theta(nT_na))(s-1)
=n\,n^{s-1}(\Theta a)(s-1)
=n^s(\Theta a)(s-1).
\tag{CSP11.13}
\] Multiplication by \(1-E_+\) commutes with \(n^s\); hence
\(\mathfrak n(nT_na)=T_n\mathfrak n(a)\). This proves the exact morphism
between the geometric normal factor and the existing translated-divisor
kernel.

The full analytic multiplier of that extension remains \[
F_0(s)F_0(s-1)=
\left[\frac{s(s-1)}8\pi^{-s/2}\Gamma(s/2)\zeta(s)\right]
\left[\frac{(s-1)(s-2)}8\pi^{-(s-1)/2}
\Gamma((s-1)/2)\zeta(s-1)\right].
\tag{CSP11.14}
\] The second factor has values \(1/8\) at \(s=1,2\), and at \(s=1-2r\)
its value is precisely (CSP11.7). Every derivative is the corresponding
derivative of \(F_0\) at \(s-1\). These are translations of the original
full formula, not deletions of its poles or trivial-zero cancellations.

On an actual original zero block of multiplicity \(m_\rho\), the normal
quotient action is \[
n^{\rho+1}
\sum_{j=0}^{m_\rho-1}\frac{(\log n)^j}{j!}M_t^j,
\qquad t=s-\rho.
\tag{CSP11.15}
\] Translation identifies this with the block at \(\rho+1\), with local
parameter \(s-(\rho+1)\). Every nilpotent coefficient, factorial and
multiplicity is retained.

\subsection{CSP12. A comparison of the whole two-degree receiver with
GSL}\label{csp12.-a-comparison-of-the-whole-two-degree-receiver-with-gsl}

The geometric calculation has produced the original quotient in degree
one and its degree-shifted coefficient in degree two. These degrees must
remain visible. Degree one is already Hausdorff by (CSP1.7); degree two
is passed through the actual quotient (CSP10.4). Write \[
\mathcal H^1=H^1(Y,\mathscr F)=H^1(Y,\mathscr F)_H\simeq\mathcal Q,\qquad
\mathcal H^2=H^2(Y,j_!\underline{Q_H})\simeq\mathcal Q(-1).
\tag{CSP12.1}
\] The combined covering action is \(T_n\) on the first and \(nT_n\) on
the second.

There is a continuous representation isomorphism of their underlying
direct sum with the GSL extension: \[
\Psi_Y:\mathcal H^1\oplus\mathcal H^2
\longrightarrow\mathcal Q_{\cup,+},
\] \[
\Psi_Y([f]_{\mathcal I},[h]_{\mathcal I})
=[E_+(s)f(s)+(1-E_+(s))h(s-1)]_{\mathcal I\cap\mathcal I_+}.
\tag{CSP12.2}
\] The second displayed input has the twisted action in (CSP12.1).
Independence of its representatives follows from the full original and
shifted jet conditions on \(E_+\). Its inverse takes a common
representative \(F\) to \[
\bigl([F]_{\mathcal I},[F(s+1)]_{\mathcal I}\bigr).
\tag{CSP12.3}
\] For the second coordinate, changing \(F\) by
\(\mathcal I\cap\mathcal I_+\) changes \(F(s+1)\) by a function in
\(\mathcal I\). On original jets, (CSP12.2) is \(f\); on shifted jets it
is \(h(s-1)\). These two observations prove both inverse identities. The
GSL global bounds prove continuity of the forward map, and translation
plus quotient seminorms prove continuity of the inverse. Formula
(CSP11.13) proves equivariance.

This is a proved comparison after taking the stated quotient and
forgetting the distinction between two cohomological degrees in the
target direct sum. It does not assert that the short exact GSL row
itself is the localization triangle of \(\mathscr F\), nor that an
unconstructed geometric boundary has vanished. The actual localization
maps are (CSP7.2)--(CSP7.6).

A diagram of the exact normal-direction maps is: \[
\begin{array}{ccccc}
J(-1)&\longrightarrow&A(-1)&\longrightarrow&Q_{\rm alg}(-1)\\
&&\big\downarrow{\nu}&&\big\downarrow\\
&&\mathcal Q(-1)&=&\mathcal Q(-1)\\
&&\big\downarrow{j_{\rm high}}\\
&&\mathcal Q_{\cup,+}.
\end{array}
\tag{CSP12.4}
\] The upper row has kernel \(J(-1)\). The middle-left vertical map has
the same kernel \(C(-1)=J(-1)\); the upper-right vertical map is an
isomorphism, its earlier comparison kernel \((C/J)(-1)\) being zero. The
bottom map is injective. All arrows intertwine the combined covering
actions with their displayed factors.

\subsection{CSP13. Exact scope of the derived geometric
statement}\label{csp13.-exact-scope-of-the-derived-geometric-statement}

The construction establishes the following concrete facts together:

\begin{enumerate}
\def\labelenumi{\arabic{enumi}.}
\tightlist
\item
  The inverse image of the actual CS coefficient sheaf on the complex
  projective-line evaluation has the fibre-product presentation
  (CSP2.2), keeping every source coordinate and all four endpoint lines.
\item
  Its global groups, costalks and localization maps are computed by
  (CSP4.2), (CSP6.2) and (CSP7.1)--(CSP7.7). The additional degree-two
  coefficient comes from annular cohomology, as also identified by
  \(R^1\pi_*\) in (CSP5.1).
\item
  The positive degree of the holomorphic cover \(z\mapsto z^n\) gives
  the exact factor \(n\). Combined with the original arithmetic
  coefficient action, this produces \(nT_n\), and therefore the
  geometric normal representation \(A(-1)\).
\item
  The full-source quotient of that normal coefficient maps by (CSP11.12)
  onto the entire existing GSL kernel. The exact comparison retains
  ordinary/Hausdorff kernels, the Mellin factor two, the original zeta,
  its full multipliers, every zero multiplicity, and the shifted
  endpoint values.
\item
  Holomorphic and antiholomorphic swaps act differently on degree two,
  exactly as in (CSP9.3). The latter is the real structure in the cited
  later CC construction.
\end{enumerate}

This is an explicit geometric source for the previously translated
coefficient character. The coefficient sheaf is infinite-dimensional and
complex, with the proved Fréchet/quotient topologies; it has not been
made into a finite-rank mixed étale sheaf on a finite-type scheme over a
finite field. Deligne's étale purity theorem is consequently not applied
to it by naming its normal factor. The statement proved here is the
construction and comparison of the actual objects and maps above.

The source dependencies are CS3 for the exact summation/Fourier source,
S1 for Mellin inversion, SSI1--SSI7 and its ESI0--ESI8 independent
derivation for the actual closed summation image strengthening S5, and
GSL4--GSL6 for its already constructed global continuous separator and
shifted quotient. CSP2--CSP10 give the new sheaf, resolution, support
and degree calculations in full. No independent norm, coordinate, or
parity is assigned to \(\tau\).
